\chapter{金融的本质，是陌生人之间的信任}
\chaptermeta{25}

\begin{ledetext}
走完了二十四章管道。现在回到楼下那家砂锅店，看看那 98 块钱现在长什么样。
\end{ledetext}

\section{回到那 98 块}

序章里你扫了个码，98 块。某个机房的硬盘上，你名下那行从 6432.10 变成 6334.10，老陈名下那行从 21847.60 变成 21945.60。

当时我说，这就是"付钱"的全部物理过程。你大概只当它是个说法。

现在把同一件事再看一遍。

\begin{bsheet}
  \bsside{你的账 · 资产}{%
    \bsrow{银行存款}{6334.10}
    \bsrowtotal{这是你的资产}{6334.10}
  }
  \bsside{银行的账 · 负债}{%
    \bsrow{你的存款}{6334.10}
    \bsrowtotal{这是它欠你的}{6334.10}
  }
\end{bsheet}

你手机里那个数字，从来不是"钱"。它是一家银行写给你的欠条。你把这张欠条转了 98 的份额给老陈，两家银行之间怎么算清，去了央行的账本上。你吃一顿饭，惊动了中央银行，这句话现在你知道它字面上是真的。

老陈拿这笔流水去还了一期经营贷。还掉的那部分本金，从这个世界上消失了：银行资产端的贷款少一块，负债端的存款同步少一块，货币总量真的变少。他剩下的钱买了一只 R2 理财，理财买了信用债，债的钱到了江苏一家做储能电芯的厂子，厂子订了一台德国的极片分切机，人民币换成欧元，欧元进了德国厂商的财资账户，去买了短期债券。

整条链现在应该都亮着。

\begin{chainflow}
\chainnode{钱是记账}\chainarrow \chainnode{记账要有人信}\chainarrow \chainnode{信任靠制度和抵押品}\chainarrow \chainnode{银行放贷时创造存款}\chainarrow \chainnode{钱约等于债}\chainarrow \chainnode{债要付息}\chainarrow \chainnode{付息需要增长}\chainarrow \chainnode{增长需要投资}\chainarrow \chainnode{投资需要给风险定价}\chainarrow \chainnode{定价错误加杠杆等于危机}\chainarrow \chainnode{危机后重置资产负债表}\chainarrow \chainnode{循环重来}
\end{chainflow}

二十四章之前，这行字对你是十二个词组。现在它是一条你能倒着走回来的路。

\section{它只干三件事}

把所有的产品、机构、术语、监管条文全部拿掉，剩下的只有三件事。

\textbf{跨时间搬钱。}把你四十岁的收入搬到七十岁去花，把你二十八岁买不起的房子搬到今天来住。存款、贷款、债券、养老金，都是这一件事的不同穿法。

\begin{egbox}{举个栗子：跨时间搬钱的标价}
老陈那 60 万经营贷，三年期，年利率 5.4\%，等额本息。月供 18091 元，三十六个月一共还回去 65 万 1 千多。

他今天多了一家店，代价是未来三年每个月固定支出 18091 元，以及 5 万 1 千多的利息。

这就是"把明天的钱搬到今天"的标价。所有金融产品的定价，剥到最后都是在给这一件事标价：\textbf{时间是有价格的。}第 3 章那句话，走到这里正好绕回来。

\end{egbox}

\textbf{在人群中分摊风险。}一个人扛不起的事，一万个人一起扛就不算事。保险、股份公司、再保险、衍生品，本质都是把一个人的悬崖，摊成一万个人的台阶。

\textbf{让陌生人之间的承诺变得可执行。}

第三件最难，也是这本书书名的那个意思。前两件事在任何有血缘和村庄的社会里都存在，只是规模小得可怜。真正把规模放大到几十亿人的，是第三件。

\section{最不性感的那些东西}

问一个具体的问题：老陈想再开一家分店，缺 60 万。

\begin{egbox}{举个栗子：银行凭什么借给一个陌生人 60 万}
信贷员根本不认识老陈。他看的是三样东西。

一，三年的收单流水和纳税记录。这依赖会计准则和支付清算系统，让"这家店到底赚不赚钱"变成一个可以被外人核验的数字。

二，有没有能抵押的东西。老陈名下那套小两居能不能抵押，取决于不动产登记系统能不能在一天之内把抵押权登记上，并且这个登记在法庭上算数。

三，还不上怎么办。这依赖破产法和强制执行：清偿顺序是可预期的，抵押物是能被处置的。

三样里缺一样，老陈的利率就从 5\% 跳到 18\%；缺两样，他只能去找亲戚。这就是第 13 章那句"风险的价格"落地之后的样子。

\end{egbox}

会计准则、清算系统、破产法、抵押登记、审计、监管。没有一样能上头条，没有一样有故事性，绝大多数人一辈子不会主动读它们中的任何一行。

\textbf{但信任的物理基础，全在这里。}它们让"我凭什么信一个我永远不会见面的人"这个问题，有了一个不依赖人品的答案。

\begin{pullquote}{}
一个国家的利率下限，很大程度上是由它的法院和登记处决定的，不是由它的央行决定的。
\end{pullquote}

所以我愿意下一个不太讨喜的判断：\textbf{一个社会金融体系的上限，不是由它最聪明的交易员决定的，是由它最不被注意的那些规则决定的。}金融创新的速度可以很快，信任的地基长得很慢。地基长不上去的地方，创新出来的东西最后都会变成第 23 章里的某一行。

\begin{mythbox}{常见误解}
\mvfrow{误解}{"金融不创造任何价值，就是一群人在互相倒手，本质是零和赌场。"}
\mvfrow{事实}{交易环节确实有大量零和的部分，这一点批评得对。但把资金从"手里有钱但没项目的人"送到"有项目但没钱的人"手里，这件事本身创造价值：老陈的第二家店雇了 7 个人，那台德国分切机让电芯厂的良品率上去了。问题不在于金融该不该存在，而在于它有多大比例真的在干这件事，多大比例只是在自己内部循环收费。这个比例，全世界都还没算清楚。}
\end{mythbox}

\bookbreak

\section{三个从来没变过的问题}

雅浦岛沉在海里的石头、罗马的金币、宋朝的交子、汇丰的银行券、你的支付宝余额、链上的稳定币。

技术换了六次，问题一次没换。

\textbf{谁在承诺？}这张欠条是谁写的。是主权国家、是持牌银行、是一家开曼注册的公司，还是一段没人负责的代码。

\textbf{凭什么信他？}他背后有什么：税收权、储备资产、抵押品、法律约束，还是仅仅是一份 PDF。

\textbf{出事了谁兜底？}存款保险、央行的最后贷款人、破产清算顺序，还是"祝你好运"。

拿这三句去问 2026 年任何一个新东西，比如稳定币：发行方是一家私人公司；它凭储备资产（短期国债）和 GENIUS Act 建立的监管框架；出事了看储备够不够、赎回条款怎么写，没有存款保险。三句话问完，你对它的位置就有数了。它不是骗局，也不是新黄金，它是一张利息归别人的、跑得很快的欠条。

\section{九条}

不凑十条。这九条是从前面二十四章里长出来的，不是我编的口号。

\begin{flowlist}
  \flowitem{01}{任何"钱"都是某个人的负债}{你的存款是银行欠你的，现金是央行欠你的，稳定币是发行方欠你的。搞清楚是谁欠你，比搞清楚有多少更重要。}
  \flowitem{02}{一个人的资产是另一个人的负债}{所以"全社会同时去杠杆"在算术上做不到，只能转移。谁接走了，是每一次危机真正的争论点。}
  \flowitem{03}{高收益不是被发现的，是被承担的}{没有藏在角落里等你捡的高收益。多出来的每一个点，都对应着某种你必须扛下来的东西。}
  \flowitem{04}{能不能随时拿回来，比收益率重要}{流动性是唯一一个不能被话术绕过的检验。}
  \flowitem{05}{不要用会被追缴的钱做期限长的事}{短钱长投是所有爆仓的同一个死法，从 1998 年的 LTCM 到你隔壁那个用信用卡套现炒股的人，机制完全一样。}
  \flowitem{06}{先活下来，再谈收益率}{任何数乘以 0 都是 0。你只有一条时间线。}
  \flowitem{07}{储蓄率是你唯一能完全控制的变量}{收益率不听你的，市场不听你的，房价不听你的。每月存多少听你的。}
  \flowitem{08}{波动不等于风险}{永久性损失才是风险。但如果波动会逼你在低点卖出，它就变成了风险，这也是为什么现金缓冲比看上去重要。}
  \flowitem{09}{你看到的收益率，都是幸存者的收益率}{包括所有的历史数据、所有的排行榜、所有的成功故事。包括这本书引用的那些。}
\end{flowlist}

\section{我们不知道的事}

一本在这些问题上装作有答案的书，两年就过期了。所以老实列在这里。

AI 究竟会不会提高全要素生产率，还是只是把成本从一个部门挪到另一个部门。这轮资本开支已经花掉了，产出还没到。

2\% 是不是最优的通胀目标。这个数字最早来自 1980 年代末新西兰的一次政策讨论，从来没有严格的理论证明。

货币政策的传导到底有多强。2021 到 2023 年的那轮通胀里，加息和通胀回落之间有多少是因果，学界还在吵。

加密资产最终会变成什么。支付基础设施、投机资产、监管套利工具，还是三样都有一点。

被动投资占比的临界点在哪里，或者到底有没有临界点。

这五个问题，2026 年 8 月没有人知道答案。任何告诉你他知道的人，你现在应该有能力判断他为什么要这么说。

\begin{takeaway}{本章一句话}
金融是一整套让陌生人之间的承诺变得可执行的技术。技术一直在换，那三个问题一次没换：谁在承诺，凭什么信他，出事了谁兜底。

\begin{itemize}
  \item 它只干三件事：跨时间搬钱、在人群中分摊风险、让陌生人的承诺可执行。
  \item 信任的地基是会计、清算、破产法和登记，不是聪明的交易员。
  \item 不做金融的信徒，也不做它的敌人。做一个看得懂它的人。
\end{itemize}

\end{takeaway}

\begin{databox}{最后一次提醒时间口径}
本书写于 2026 年 8 月。书里所有的数字都带着这个日期：美国联邦基金利率 3.50\%–3.75\%，中国 5 年期以上 LPR 3.5\%，全球央行储备里黄金占 27\%，全球债务约 353 万亿美元。

这些数字明年就不对了。管道明年还在。

\end{databox}

上个月老陈把价格调了。两个菜一碗汤，现在 104。

他说是猪骨和菜价都涨了。他不知道的是，这 6 块钱里有一小块是中东的油价，还有一小块是他去年那笔经营贷的利息。

你知道。

汤还是好喝的。

\begin{bridgebox}
\textbf{下一站}附录：术语速查表，以及 2026 年 8 月的数据快照。前者用来查，后者用来两年后回头对照，看看哪些变了、哪些没变。
\end{bridgebox}

